

\section{Reduction} \label{section:Reduction}

In this section, we give a procedure to reduce the proof of Theorem \ref{thm:MainTarget} to a classification problem.

\subsection{Setting and conventions}

From this section onward, we use the following convention.
For a real reductive Lie algebra $\lie{g}$ with a Cartan involution $\theta$, we fix a connected reductive Lie group $G$ with the Lie algebra $\lie{g}$ such that $G$ has a Cartan involution $\theta$ whose differential is the Cartan involution $\theta$ on $\lie{g}$.
Fix also a $G$-invariant non-degenerate symmetric bilinear form $(\cdot, \cdot)$ on $\lie{g}$ such that $-(\cdot, \theta(\cdot))$ is positive definite and symmetric.
Extend $(\cdot, \cdot)$ to a bilinear form on $\lieC{g}$.
Set $K \coloneq G^\theta$.
For a $\theta$-stable subgroup $H$ of $G$, we set $K_H \coloneq H \cap K$.

Throughout this section, for a Lie subalgebra of $\lie{g}$, we denote by the corresponding capital Roman letter the analytic subgroup of $G$ with that Lie algebra; for example, $H$ denotes the analytic subgroup of $G$ with Lie algebra $\lie{h}$.
In several propositions, we use reductions involving analytic subgroups.
We neither assume nor verify that these analytic subgroups are closed.
The arguments are valid without the closedness assumption.

We say that a $\theta$-stable subalgebra $\lie{h}$ of $\lie{g}$ is compact if $\lie{h} \subset \lie{k}$.
Note that $H$ may not be compact even if $\lie{h}$ is compact because $H$ may not be closed in $G$.
Clearly, $H$ is relatively compact in $K$.

\subsection{Minimal tuples}

In this subsection, we prepare several notions and lemmas to reduce the proof of Theorem \ref{thm:MainTarget} to small cases.
Let $\lie{g}$ be a real reductive Lie algebra with Cartan involution $\theta$, $\lie{h}$ a $\theta$-stable reductive subalgebra of $\lie{g}$ and $\orbit{O}$ a nilpotent orbit in $\lie{g}$.

\begin{definition}
	We say that the triple $(\lie{g}, \lie{h}, \orbit{O})$ is NDD (non-discretely decomposable) if $\overline{\orbit{O}} \cap \Nilpotent(\lie{h}^\perp, H) \neq \set{0}$.
\end{definition}

\begin{definition} \label{def:MainAssertion}
	We denote by $\assertion(\lie{g}, \lie{h}, \orbit{O})$ the assertion that $\proj_{\lieC{h}}(\KS(\overline{\orbit{O}})) \not \subset \Nilpotent((\lie{k}_{H, \CC})^\perp, K_{H, \CC})$.
\end{definition}

\begin{remark}
	The two definitions are independent of the choice of the bilinear form $(\cdot, \cdot)$ as follows.
	In fact, the following diagram commutes:
	\begin{align*}
		\xymatrix{
			\lieC{g} \ar[r] \ar[d]^{\proj_{\lieC{h}}} & \lieC{g}^* \ar[d]^{\mathrm{res}_{\lieC{h}}} & \lieC{g} \ar[l] \ar[d]^{\proj'_{\lieC{h}}} \\
			\lieC{h} \ar[r] & \lieC{h}^* & \lieC{h} \ar[l],
		}
	\end{align*}
	where $\proj'_{\lieC{h}}$ is the orthogonal projection with respect to another $G$-invariant non-degenerate symmetric bilinear form, the horizontal maps are linear isomorphisms induced by the bilinear forms and $\mathrm{res}_{\lieC{h}}$ is the restriction map to $\lieC{h}$.
	Clearly, all the maps in the diagram are $\lieC{h}$-equivariant.
	Moreover, since $\orbit{O}$ is closed under the $\RR_{>0}$-action on each simple ideal of $\lie{g}$, the bijection $\Nilpotent(\lie{g}, G)/G \simeq \Nilpotent(\lie{g}^*, G)/G$ is independent of the choice of the bilinear form.
\end{remark}

It is obvious that Theorem \ref{thm:MainTarget} follows if the assertion $\assertion(\lie{g}, \lie{h}, \orbit{O})$ is proved for any NDD triple $(\lie{g}, \lie{h}, \orbit{O})$.
Hereafter, assume that $(\lie{g}, \lie{h}, \orbit{O})$ is NDD.
Let $\lie{a}_H$ be a maximal abelian subspace of $\lie{h}^{-\theta}$.

\begin{lemma}\label{lem:ExistenceWeightVector}
	Let $\lie{a}$ be a maximal abelian subspace of $\lieNorm_{\lie{g}}(\lie{h})^{-\theta}$ containing $\lie{a}_H$.
	There exists an $\lie{a}$-weight vector $X \in \overline{\orbit{O}} \cap \Nilpotent(\lie{h}^\perp, H)$ with non-zero weight on $\lie{a}_H$.
\end{lemma}

\begin{proof}
	Let $\lie{h}_{\mathrm{nc}}$ be the subalgebra of $\lie{h}$ generated by $\lie{h}^{-\theta}$.
	Assume that any element in $\overline{\orbit{O}} \cap \Nilpotent(\lie{h}^\perp, H)$ is $\lie{a}_H$-invariant.
	Then any element in $\overline{\orbit{O}} \cap \Nilpotent(\lie{h}^\perp, H)$ is $\lie{h}_{\mathrm{nc}}$-invariant since $\overline{\orbit{O}}\cap \Nilpotent(\lie{h}^\perp, H)$ is $K_H$-stable.
	For any $X \in \overline{\orbit{O}} \cap \Nilpotent(\lie{h}^\perp, H)$, we have
	\begin{align*}
		0 \in \overline{\Ad(H)X} \subset \set{Z \in \lie{g} : (Z, \theta(Z)) = (X, \theta(X))}
	\end{align*}
	and hence $X = 0$.
	This contradicts the assumption that $(\lie{g}, \lie{h}, \orbit{O})$ is NDD.

	Thus, there exists $X \in \overline{\orbit{O}} \cap \Nilpotent(\lie{h}^\perp, H)$ that is not $\lie{a}_H$-invariant.
	Write $X = X_1 + X_2 + \cdots + X_k$ according to the $\lie{a}_H$-weight space decomposition, where $X_i$ has weight $\lambda_i$.
	The convex hull of $\set{\lambda_1, \lambda_2, \ldots, \lambda_k}$ is not $\set{0}$.
	Hence there exist $\lambda_i$ and $H \in \lie{a}_H$ such that $\lambda_i \neq 0$ and $\lambda_i(H) > \lambda_j(H)$ for all $j \neq i$.
	Since $\overline{\orbit{O}}$ is a $G$-stable closed cone, we have
	\begin{align*}
		X_i = \lim_{t \to \infty} e^{-t\lambda_i(H)}\Ad(e^{tH})X \in \overline{\orbit{O}} \cap \Nilpotent(\lie{h}^\perp, H).
	\end{align*}
	$X_i$ is the desired $\lie{a}_H$-weight vector with non-zero weight $\lambda_i$.

	Since $\lie{a}$ normalizes $\lie{h}^\perp$, we can apply the same argument to the $\lie{a}_H$-weight vector $X_i$, with $\lie{a}$ in place of $\lie{a}_H$.
	Then we obtain an $\lie{a}$-weight vector in $\overline{\orbit{O}} \cap \Nilpotent(\lie{h}^\perp, H)$ with non-zero weight on $\lie{a}_H$.
	This proves the lemma.
\end{proof}

Before discussing reductions for NDD tuples, we prepare a lemma which allows us to remove the unnecessary factors of $\lie{g}$ and $\lie{h}$.

\begin{lemma} \label{lem:ReductionIdeal}
	Let $\lie{g'}$ be a $\theta$-stable ideal of $\lie{g}$ such that $\orbit{O}\subset \lie{g'}$.
	Let $\lie{h'}$ be the orthogonal projection of $\lie{h}$ to $\lie{g'}$.
	Then $(\lie{g'}, \lie{h'}, \orbit{O})$ is NDD.
	Moreover, $\assertion(\lie{g'}, \lie{h'}, \orbit{O})$ holds if and only if $\assertion(\lie{g}, \lie{h}, \orbit{O})$ holds.
\end{lemma}

\begin{remark}
	The assumption on $\lie{g'}$ is satisfied if $(\lie{g'})^\perp$ is a direct sum of abelian ideals and compact ideals.
\end{remark}

\begin{proof}
	Clearly, we have $\lie{h}^\perp \cap \lie{g'} = (\lie{h'})^\perp \cap \lie{g'}$
	and the $H'$-orbits in $\lie{g'}$ are the same as the $H$-orbits in $\lie{g'}$.
	In particular, $(\lie{g'}, \lie{h'}, \orbit{O})$ is NDD.

	Let $Z \in \lieC{g}'$.
	It is sufficient to show that $\proj_{\lieC{h}'}(Z)$ is nilpotent if and only if $\proj_{\lieC{h}}(Z)$ is nilpotent (see Remark \ref{remark:AssociatedVariety}).
	If $Z \in \lieC{g}'\cap (\lieC{h}')^\perp$, then we have $\proj_{\lieC{h}'}(Z) = \proj_{\lieC{h}}(Z) = 0$.
	Hence we may assume that $Z \in \lieC{h}'$.
	Note that $\proj_{\lieC{h}}|_{\lieC{h}'}$ is injective and $\lieC{h}$-equivariant.
	Hence $\overline{\Ad(H'_\CC)Z} \ni 0$ if and only if $\overline{\Ad(H_\CC)\proj_{\lieC{h}}(Z)} \ni 0$.
	This implies that $\proj_{\lieC{h}}(Z)$ is nilpotent if and only if $\proj_{\lieC{h}'}(Z) = Z$ is nilpotent.
\end{proof}


\begin{lemma} \label{lem:ReductionCompactFactor}
	Let $\lie{k'}$ be a subalgebra of $\lieCent_{\lie{k}}(\lie{h}) \cap \lie{h}^\perp$.
	Then $(\lie{g}, \lie{h}, \orbit{O})$ is NDD if and only if $(\lie{g}, \lie{h} \oplus \lie{k'}, \orbit{O})$ is NDD.
	Moreover, $\assertion(\lie{g}, \lie{h}, \orbit{O})$ holds if and only if $\assertion(\lie{g}, \lie{h} \oplus \lie{k'}, \orbit{O})$ holds.
\end{lemma}

\begin{proof}
	Let $p \colon \lie{g} \rightarrow \lie{k'}$ be the orthogonal projection.
	Since $\lie{k'}$ is contained in $\lie{g}^H$, we have $p(\Nilpotent(\lie{h}^\perp, H)) \subset \Nilpotent(\lie{k'}, H) = \set{0}$, and hence
	\begin{align*}
		\overline{\orbit{O}} \cap \Nilpotent(\lie{h}^\perp, H) = \overline{\orbit{O}} \cap \Nilpotent((\lie{h} \oplus \lie{k'})^\perp, H).
	\end{align*}
	Since $K' \subset K$, the analytic subgroup $K' \subset G$ preserves the inner product $-(\cdot, \theta(\cdot))$ on $\lie{g}$.
	This implies that 
	\begin{align*}
		\overline{\orbit{O}} \cap \Nilpotent(\lie{h}^\perp, H) = \overline{\orbit{O}} \cap \Nilpotent((\lie{h} \oplus \lie{k'})^\perp, H) = \overline{\orbit{O}} \cap \Nilpotent((\lie{h} \oplus \lie{k'})^\perp, H\times K').
	\end{align*}
	Hence $(\lie{g}, \lie{h}, \orbit{O})$ is NDD if and only if $(\lie{g}, \lie{h} \oplus \lie{k'}, \orbit{O})$ is NDD.

	Recall that $\KS(\overline{\orbit{O}})$ is contained in $(\lieC{g})^{-\theta}$.
	See Fact \ref{fact:KostantSekiguchi}.
	Hence, for fixed $\lie{g}$ and $\orbit{O}$, the assertion $\assertion(\lie{g}, \lie{h}, \orbit{O})$ depends on $\lie{h}$ only through $(\lieC{h})^{-\theta}$.
	This shows that $\assertion(\lie{g}, \lie{h}, \orbit{O})$ holds if and only if $\assertion(\lie{g}, \lie{h} \oplus \lie{k'}, \orbit{O})$ holds.
\end{proof}

\begin{lemma} \label{lem:ReductionNilpotent}
	Let $X$ be an $\lie{a}_H$-weight vector in $\overline{\orbit{O}} \cap \Nilpotent(\lie{h}^\perp, H)$ with non-zero weight.
	Define $\lie{g'}$ to be the smallest subalgebra of $\lie{g}$ that contains $\lie{a}_H$, $X$ and $\theta(X)$, and satisfies $\lie{g'} = (\lie{g'}\cap \lie{h}) \oplus (\lie{g'}\cap \lie{h}^\perp)$.
	Set $\lie{h'} \coloneq \lie{g'} \cap \lie{h}$ and $\orbit{O}' \coloneq \Ad(G')X$.
	Then $(\lie{g'}, \lie{h'}, \orbit{O}')$ is NDD.
	Moreover, if $\assertion(\lie{g}', \lie{h}', \orbit{O}')$ holds, then so does $\assertion(\lie{g}, \lie{h}, \orbit{O})$.
\end{lemma}

\begin{proof}
	By construction, the weight vector $X$ belongs to $\orbit{O}'\cap \Nilpotent(\lie{g'}\cap \lie{h}^\perp, H')$.
	Hence $(\lie{g'}, \lie{h'}, \orbit{O}')$ is NDD.
	
	Since $\lie{g'} = (\lie{g'}\cap \lie{h}) \oplus (\lie{g'}\cap \lie{h}^\perp)$, we have $\proj_{\lieC{h}}|_{\lieC{g}'} = \proj_{\lieC{h'}}|_{\lieC{g}'}$. Recall that $\KS(\overline{\orbit{O}'}) \subset \KS(\overline{\orbit{O}})$ by Proposition \ref{prop:KostantSekiguchiCompatible}.
	The second assertion follows from this.
\end{proof}

In Lemmas \ref{lem:ReductionIdeal}, \ref{lem:ReductionCompactFactor} and \ref{lem:ReductionNilpotent}, we have given reductions for NDD tuples.
We shall state them explicitly.

\begin{definition}\label{def:Reduction}
	We define three reductions for NDD tuples $(\lie{g}, \lie{h}, \orbit{O})$ as follows.
	\begin{enumerate}[label=(\alph*)]
		\item Let $\lie{g}'$ be the smallest ideal of $\lie{g}$ containing $\orbit{O}$.
		Define $\lie{h}'$ as in Lemma \ref{lem:ReductionIdeal}.
		Replace $(\lie{g}, \lie{h}, \orbit{O})$ by $(\lie{g'}, \lie{h'}, \orbit{O})$.
		\item Let $\lie{i}$ be the sum of all compact ideals of $\lie{h}$ and $\lie{h'}$ be the orthogonal complement of $\lie{i}$ in $\lie{h}$.
		Replace $(\lie{g}, \lie{h}, \orbit{O})$ by $(\lie{g}, \lie{h'}, \orbit{O})$.
		\item Choose a $G$-invariant non-degenerate symmetric bilinear form $(\cdot, \cdot)$ on $\lie{g}$ such that $-(\cdot, \theta(\cdot))$ is positive definite and symmetric.
		Let $X$ be an $\lie{a}_H$-weight vector in $\overline{\orbit{O}} \cap \Nilpotent(\lie{h}^\perp, H)$ with non-zero weight.
		Define $\lie{g'}$ and $\lie{h'}$ as in Lemma \ref{lem:ReductionNilpotent}.
		Replace $(\lie{g}, \lie{h}, \orbit{O})$ by $(\lie{g'}, \lie{h'}, \Ad(G')X)$.
	\end{enumerate}
\end{definition}

\begin{remark} \label{remark:minimal}
	In Definition \ref{def:Reduction} (c), the choice of the $G$-invariant bilinear form $(\cdot, \cdot)$ on $\lie{g}$ is crucial.
	Let $\lie{h} \coloneq \lieSL(2,\RR) \hookrightarrow \lieSL(2,\RR)\oplus \lieSL(2,\RR) \eqcolon \lie{g}$ be the diagonal embedding.
	Consider the Killing form on $\lie{g}$.
	Then $\lie{h}^\perp$ is the anti-diagonal embedding of $\lieSL(2,\RR)$, and hence $\lie{h}^\perp$ contains a non-zero nilpotent element $X$ such that $\set{[X, -\theta(X)], X, -\theta(X)}$ forms an $\lieSL_2$-triple.
	In other words, $\lie{g}'$ is a proper subalgebra of $\lie{g}$ for the choice of $X$.
	
	On the other hand, consider the bilinear form on $\lie{g}$ defined by the direct sum of the Killing form on the first factor and some scalar multiple of the Killing form on the second factor.
	Then $\lie{h}^\perp$ may not contain any non-zero nilpotent element $X$ such that $\set{[X, -\theta(X)], X, -\theta(X)}$ forms an $\lieSL_2$-triple.
	Hence $\lie{g}' = \lie{g}$ for any choice of $X$.
\end{remark}

\begin{definition} \label{def:minimal}
	If $(\lie{g}, \lie{h}, \orbit{O})$ is unchanged under every choice of the reductions (a), (b) and (c) in Definition \ref{def:Reduction}, we say that the tuple $(\lie{g}, \lie{h}, \orbit{O})$ is minimal.
\end{definition}

For any NDD tuple $(\lie{g}, \lie{h}, \orbit{O})$, after finitely many reductions (a), (b) and (c) in Definition \ref{def:Reduction}, one obtains a minimal tuple.
Note that there are only finitely many nilpotent orbits in $\lie{g}$.

The minimality condition is so strong that essentially only one case can occur.
By Lemmas \ref{lem:ReductionIdeal}, \ref{lem:ReductionCompactFactor} and \ref{lem:ReductionNilpotent}, our goal is reduced to proving the following theorem.

\begin{theorem} \label{thm:MainTargetMinimal}
	If an NDD tuple $(\lie{g}, \lie{h}, \orbit{O})$ is minimal, then one has $(\lie{g}, \lie{h}) \simeq (\lieSL(2,\RR), \RR)$, where $\RR$ means the subalgebra of diagonal matrices in $\lieSL(2,\RR)$.
\end{theorem}

\begin{remark}
	If $(\lie{g}, \lie{h}) \simeq (\lieSL(2,\RR), \RR)$, the assertion $\assertion(\lie{g}, \lie{h}, \orbit{O})$ is easy to verify.
	It is also easy to see that $(\lie{g}, \lie{h}, \orbit{O})$ is minimal.
\end{remark}

\begin{remark}
	In the proof of \cite[Theorem 21]{Ki24}, we used the existence of an $\lieSL_2$-triple $\set{[X, -\theta(X)], X, -\theta(X)}$ such that $X \in \Nilpotent(\lie{h}^\perp, H) \cap \overline{\orbit{O}}$.
	Theorem \ref{thm:MainTargetMinimal} does not assert that, for any NDD tuple $(\lie{g}, \lie{h}, \orbit{O})$, there exists such an $\lieSL_2$-triple.
	In fact, this is not true when $(\lie{g}, \lie{h}) = (\lieSO(4, 3), \lieSU(2, 1))$ and $(\lie{g}, \lie{h}) = (\lieG_{2(2)}, \lieSU(2, 1))$.
\end{remark}

If $\lie{h}$ is abelian, Theorem \ref{thm:MainTargetMinimal} is easy to prove.
We will see this in Theorem \ref{thm:MinimalAbelian}.
The remaining proof of Theorem \ref{thm:MainTargetMinimal} will be given in Sections \ref{section:TypeBC1} and \ref{section:TypeA1}.
In the remaining part of this section, we give several necessary conditions for minimal tuples.

To show the non-minimality of an NDD tuple $(\lie{g}, \lie{h}, \orbit{O})$, it is enough to find a suitable subalgebra $\lie{g}'$ of $\lie{g}$ instead of finding a suitable nilpotent element $X$ in $\lie{h}^\perp \cap \overline{\orbit{O}}$.

\begin{proposition} \label{prop:NonMinimalSubalgebra}
	Assume that there exists a $\theta$-stable reductive subalgebra $\lie{g}'\subsetneq \lie{g}$ satisfying the following conditions.
	\begin{enumparen}
		\item $\lie{g}'$ is $\lie{a}_H$-stable.
		\item $\lie{g}' = (\lie{g}' \cap \lie{h}) \oplus (\lie{g}' \cap \lie{h}^\perp)$.
		\item There exists an $\lie{a}_H$-weight vector $X$ in $\lie{g}' \cap \lie{h}^\perp \cap \overline{\orbit{O}}$ with non-zero weight.
	\end{enumparen}
	Then $(\lie{g}, \lie{h}, \orbit{O})$ is not minimal.
\end{proposition}

\begin{remark}
	An $\lie{a}_H$-weight vector in $\lie{h}^\perp$ with non-zero weight automatically belongs to $\Nilpotent(\lie{h}^\perp, H)$.
\end{remark}

\begin{proof}
	If $\lie{g}$ has non-trivial center, then $(\lie{g}, \lie{h}, \orbit{O})$ is not minimal by definition.
	Hence we may assume that $\lie{g}$ is semisimple.

	Since $\lie{g}'$ is $\lie{a}_H$-stable, $\lie{g}' + \lie{a}_H$ is a subalgebra of $\lie{g}$.
	Let $\lie{g}''$ be the smallest subalgebra of $\lie{g}$ that contains $\lie{a}_H$, $X$ and $\theta(X)$, and satisfies $\lie{g}'' = (\lie{g}''\cap \lie{h}) \oplus (\lie{g}''\cap \lie{h}^\perp)$.
	By assumption, we have $\lie{g}'' \subset \lie{g}' + \lie{a}_H$.
	This implies that $[\lie{g}'', \lie{g}''] \subset \lie{g}' \subsetneq \lie{g}$.
	Since $\lie{g}$ is semisimple, $\lie{g}''$ is a proper subalgebra of $\lie{g}$.
	This means that reduction (c) in Definition \ref{def:Reduction} replaces $\lie{g}$ by the proper subalgebra $\lie{g}''$.
	This shows that $(\lie{g}, \lie{h}, \orbit{O})$ is not minimal.
\end{proof}

\begin{corollary} \label{cor:MinimalGeneratedBy}
	Let $X$ be an $\lie{a}_H$-weight vector in $\overline{\orbit{O}} \cap \Nilpotent(\lie{h}^\perp, H)$ with non-zero weight $\lambda$.
	Assume that $(\lie{g}, \lie{h}, \orbit{O})$ is minimal.
	Then each $\lie{a}_H$-weight in $\lie{g}$ is of the form $n\lambda$ ($n \in \ZZ$).
\end{corollary}

\begin{proof}
	Set $\lie{g'} \coloneq \bigoplus_{n \in \ZZ} \lie{g}_{n\lambda}$.
	Then we have $X, \theta(X) \in \lie{g'}$ and $\lie{a}_H \subset \lie{g'}$.
	Since $\lie{g'}$ is $\theta$-stable and satisfies the assumptions in Proposition \ref{prop:NonMinimalSubalgebra}, we have $\lie{g'} =\lie{g}$ by minimality.
	This shows the assertion.
\end{proof}

The following proposition follows from the definition of minimal tuples and Corollary \ref{cor:MinimalGeneratedBy}.

\begin{proposition} \label{prop:BasicMinimal}
	Assume that $(\lie{g}, \lie{h}, \orbit{O})$ is minimal.
	Then the following conditions hold.
	\begin{enumparen}
		\item $\lie{g}$ is semisimple without any compact factor.
		\item $\lie{h}$ has no compact factor.
		\item $\lie{h}$ has real rank one.
		\item $\lie{h}$ is simple or $1$-dimensional non-compact abelian.
		\item $\orbit{O} = \Ad(G)X$ for any non-zero $X \in \overline{\orbit{O}} \cap \Nilpotent(\lie{h}^\perp, H)$.
		\item The projection of any non-zero element in $\overline{\orbit{O}} \cap \Nilpotent(\lie{h}^\perp, H)$ to each simple factor of $\lie{g}$ is non-zero.
	\end{enumparen}
\end{proposition}

\begin{proposition} \label{prop:BasicMinimalH}
	Assume that $(\lie{g}, \lie{h}, \orbit{O})$ is minimal.
	Then the projection of $\lie{h}$ to each simple factor of $\lie{g}$ is non-zero.
\end{proposition}

\begin{proof}
	Let $\lie{i}$ be a simple ideal of $\lie{g}$ and $p\colon \lie{g} \rightarrow \lie{i}$ be the projection.
	Take a non-zero $X \in \overline{\orbit{O}} \cap \Nilpotent(\lie{h}^\perp, H)$.
	By Proposition \ref{prop:BasicMinimal} (6), we have
	\begin{align*}
		0 \neq p(X) \in p(\Nilpotent(\lie{h}^\perp, H)) \subset \Nilpotent(\lie{i}, H).
	\end{align*}
	This implies that $p(\lie{h})$ is non-zero.
\end{proof}

\subsection{Necessary conditions for minimality}

In \cite[Theorem 21]{Ki24}, we have given sufficient conditions for the implication (3) $\Rightarrow$ (1) in Fact \ref{fact:Implication}.
In this subsection, we rewrite the sufficient conditions in the context of NDD tuples.
Let $(\lie{g}, \lie{h}, \orbit{O})$ be an NDD tuple.
Assume that $\lie{g} \not \simeq \lieSL(2,\RR)$.

\begin{lemma} \label{lem:NonMinimalX}
	Assume that there exists an $\lie{a}_H$-weight vector $X \in \lie{h}^\perp \cap \overline{\orbit{O}}$ with non-zero weight such that $\set{[X, -\theta(X)], X, -\theta(X)}$ spans a subalgebra of $\lie{g}$ isomorphic to $\lieSL(2,\RR)$.
	Then $(\lie{g}, \lie{h}, \orbit{O})$ is not minimal.
\end{lemma}

\begin{proof}
	To show the assertion, we may assume that $\lie{g}$ is semisimple.
	Let $\lie{g}'$ be the subalgebra generated by $X, -\theta(X)$ and $\lie{a}_H$.
	By assumption, $[\lie{g}', \lie{g}']$ is isomorphic to $\lieSL(2,\RR)$ and hence $\lie{g'}$ is a proper subalgebra of $\lie{g}$.
	It is easy to see that $\lie{g}'$ satisfies the assumptions in Proposition \ref{prop:NonMinimalSubalgebra}.
	Hence $(\lie{g}, \lie{h}, \orbit{O})$ is not minimal.
\end{proof}

\begin{lemma} \label{lem:SufficientSymmetric}
	Let $\lie{k}'$ be a subalgebra of $\lieCent_{\lie{k}}(\lie{h}) \cap \lie{h}^\perp$.
	If $(\lie{g}, \lie{h} \oplus \lie{k}')$ is a symmetric pair, then $(\lie{g}, \lie{h}, \orbit{O})$ is not minimal.
\end{lemma}

\begin{proof}
	Take an involution $\sigma$ of $\lie{g}$ commuting with $\theta$ such that $\lie{g}^\sigma = \lie{h} \oplus \lie{k}'$.
	If necessary, replacing the form $(\cdot, \cdot)$ on $\lie{g}$ by $(\cdot, \cdot) + (\sigma(\cdot), \sigma(\cdot))$, we may assume that $(\lie{h} \oplus \lie{k}')^\perp = \lie{g}^{-\sigma}$.

	By Lemma \ref{lem:ExistenceWeightVector}, there exists an $\lie{a}_H$-weight vector $X \in \overline{\orbit{O}} \cap \Nilpotent(\lie{h}^\perp, H)$ with non-zero weight.
	Note that $X \in (\lie{h}\oplus \lie{k}')^\perp$ as in the proof of Lemma \ref{lem:ReductionCompactFactor}.
	Since $(\lie{g}, \lie{h} \oplus \lie{k}')$ is a symmetric pair, we have $[X, -\theta(X)] \in (\lie{h} \oplus \lie{k}') \cap \lie{g}^{-\theta} \subset \lie{h}$.
	Moreover, since $X$ is a weight vector, $[X, -\theta(X)]$ commutes with $\lie{a}_H$.
	This implies $[X, -\theta(X)] \in \lie{a}_H$.
	Hence $\set{[X, -\theta(X)], X, -\theta(X)}$ spans a subalgebra of $\lie{g}$ isomorphic to $\lieSL(2,\RR)$.
	By Lemma \ref{lem:NonMinimalX}, the assertion follows.
\end{proof}

\begin{lemma} \label{lem:SufficientMaxAbelian}
	Let $\lie{a}$ be a maximal abelian subspace of $\lie{g}^{-\theta}$ containing $\lie{a}_H$.
	Assume that there exists an $\lie{a}$-weight vector $X$ in $\overline{\orbit{O}} \cap \Nilpotent(\lie{h}^\perp, H)$ with non-zero weight on $\lie{a}_H$.
	Then $(\lie{g}, \lie{h}, \orbit{O})$ is not minimal.
\end{lemma}

\begin{proof}
	By a standard argument, $\set{[X, -\theta(X)], X, -\theta(X)}$ spans a subalgebra of $\lie{g}$ isomorphic to $\lieSL(2,\RR)$.
	Lemma \ref{lem:NonMinimalX} shows that $(\lie{g}, \lie{h}, \orbit{O})$ is not minimal.
\end{proof}

\begin{corollary} \label{cor:SufficientSameRank}
	If $\lieNorm_{\lie{g}}(\lie{h})$ has the same real rank as $\lie{g}$, then $(\lie{g}, \lie{h}, \orbit{O})$ is not minimal.
\end{corollary}

\begin{proof}
	Take a maximal abelian subspace $\lie{a}$ of $\lieNorm_{\lie{g}}(\lie{h})^{-\theta}$ containing $\lie{a}_H$.
	Then $\lie{a}$ is a maximal abelian subspace of $\lie{g}^{-\theta}$ by assumption.
	By Lemma \ref{lem:ExistenceWeightVector}, there exists an $\lie{a}$-weight vector $X \in \overline{\orbit{O}} \cap \Nilpotent(\lie{h}^\perp, H)$ with non-zero weight on $\lie{a}_H$.
	Hence the assertion follows from Lemma \ref{lem:SufficientMaxAbelian}.
\end{proof}

\begin{corollary} \label{cor:SufficientNonRoot}
	Let $X$ be an $\lie{a}_H$-weight vector in $\overline{\orbit{O}} \cap \Nilpotent(\lie{h}^\perp, H)$ with non-zero weight $\lambda$.
	Assume that the weight space $\lie{g}_\lambda$ is contained in $\lie{h}^\perp$.
	Then $(\lie{g}, \lie{h}, \orbit{O})$ is not minimal.
\end{corollary}

\begin{proof}
	Take a maximal abelian subspace $\lie{a}$ of $\lie{g}^{-\theta}$ containing $\lie{a}_H$.
	By assumption, the weight space $\lie{g}_\lambda$ is contained in $\lie{h}^\perp$ and $\lie{a}$-stable.
	Hence, by the same argument as in the proof of Lemma \ref{lem:ExistenceWeightVector}, there exists an $\lie{a}$-weight vector $X' \in \lie{g}_\lambda \cap \overline{\orbit{O}} \cap \Nilpotent(\lie{h}^\perp, H)$.
	The assertion follows from Lemma \ref{lem:SufficientMaxAbelian}.
\end{proof}

As a consequence of Corollary \ref{cor:SufficientSameRank}, we prove Theorem \ref{thm:MainTargetMinimal} when $\lie{h}$ is abelian.

\begin{theorem} \label{thm:MinimalAbelian}
	Let $(\lie{g}, \lie{h}, \orbit{O})$ be a minimal NDD tuple.
	Assume that $\lie{h}$ is abelian.
	Then $(\lie{g}, \lie{h}) \simeq (\lieSL(2,\RR), \RR)$, where $\RR$ means the subalgebra of diagonal matrices in $\lieSL(2,\RR)$.
\end{theorem}

\begin{proof}
	By Proposition \ref{prop:BasicMinimal} (4), $\lie{h}$ is non-compact.
	Hence $\lieNorm_{\lie{g}}(\lie{h})$ contains a maximal abelian subspace of $\lie{g}^{-\theta}$.
	By Corollary \ref{cor:SufficientSameRank}, we have $\lie{g} \simeq \lieSL(2,\RR)$.
	Therefore, we obtain $(\lie{g}, \lie{h}) \simeq (\lieSL(2,\RR), \RR)$.
\end{proof}

\subsection{Reduction to rank one subalgebras} \label{subsect:ReductionRankOne}

In this subsection, we shall see that the minimality of an NDD tuple $(\lie{g}, \lie{h}, \orbit{O})$ strongly restricts the structure of $\lie{g}$ and $\lie{h}$.

\begin{lemma}
	Let $(\lie{g}, \lie{h}, \orbit{O})$ be a minimal NDD tuple.
	Then $\lieNorm_{\lie{g}}(\lie{h})$ has real rank one.
	In particular, $\lieCent_{\lie{g}}(\lie{h})$ is compact if $\lie{h}$ is simple.
\end{lemma}

\begin{proof}
	Assume that $\lieNorm_{\lie{g}}(\lie{h})$ has real rank greater than one.
	Take a maximal abelian subspace $\lie{a}$ of $\lieNorm_{\lie{g}}(\lie{h})^{-\theta}$ containing $\lie{a}_H$.
	By Lemma \ref{lem:ExistenceWeightVector}, there exists an $\lie{a}$-weight vector $X \in \overline{\orbit{O}} \cap \Nilpotent(\lie{h}^\perp, H)$ with weight $\lambda$ that is non-zero on $\lie{a}_H$.
	Let $\lie{g'}$ be the smallest subalgebra of $\lie{g}$ that contains $\lie{a}$, $X$ and $\theta(X)$, and satisfies $\lie{g'} = (\lie{g'}\cap \lie{h}) \oplus (\lie{g'}\cap \lie{h}^\perp)$.
	Since $\dim(\lie{a}) > 1$, we have $0 \neq \Ker(\lambda) \subset \lieCent(\lie{g'})$.
	Since $\lie{g}$ has no center, $\lie{g'}$ is a proper subalgebra of $\lie{g}$.
	By Proposition \ref{prop:NonMinimalSubalgebra}, $(\lie{g}, \lie{h}, \orbit{O})$ is not minimal.
	This is a contradiction, and we have shown the assertion.
\end{proof}

In the following, we consider a grading of $\lie{g}$ associated with an element $h \in \lie{a}_H$.
For $i \in \RR$ and $\ad(h)$-stable subspace $W \subset \lieC{g}$, we denote by $W_i$ the eigenspace of $\ad(h)$ with eigenvalue $i$.
Then we have a grading $\lie{g} = \bigoplus_{i \in \RR} \lie{g}_i$.

\begin{definition} \label{def:typeA}
	Let $\lie{g}$ be a semisimple Lie algebra and $\lie{h}$ a $\theta$-stable subalgebra.
	We say that $(\lie{g}, \lie{h})$ is of type $A_1$ if the pair satisfies the following conditions:
	\begin{enumparen}
		\item $\lie{h}$ is simple and has restricted root system of type $A_1$,
		\item $\lieCent_{\lie{g}}(\lie{h})$ is compact,
		\item There exists a non-zero element $h \in \lie{h}^{-\theta}$ such that $\lie{g} = \lie{g}_{-1} \oplus \lie{g}_0 \oplus \lie{g}_1$, and 
		\item The projection of $\lie{g}_1$ to each simple factor of $\lie{g}$ is non-zero.
	\end{enumparen}
\end{definition}

\begin{remark}
	If $(\lie{g}, \lie{h})$ is of type $A_1$, then the projection of $\lie{h}$ to each simple factor of $\lie{g}$ is non-zero.
	This follows from the fact that, by condition (4), the projection of $h$ in condition (3) to each simple factor of $\lie{g}$ is non-zero.
\end{remark}

\begin{definition} \label{def:typeBC}
	Let $\lie{g}$ be a semisimple Lie algebra and $\lie{h}$ a $\theta$-stable subalgebra.
	We say that $(\lie{g}, \lie{h})$ is of type $BC'_1$ if the pair satisfies the following conditions.
	\begin{enumparen}
		\item $\lie{g}$ is simple.
		\item $\lie{h}$ is simple and has restricted root system of type $BC_1$.
		\item $\lieCent_{\lie{g}}(\lie{h})$ is compact.
		\item There exists a non-zero element $h \in \lie{h}^{-\theta}$ such that $\lie{g} = \lie{g}_{-2} \oplus \lie{g}_{-1} \oplus \lie{g}_0 \oplus \lie{g}_1 \oplus \lie{g}_2$.
	\end{enumparen}
	We say that $(\lie{g}, \lie{h})$ is of type $BC_1$ if it is of type $BC'_1$ and satisfies the following additional condition:
	\begin{enumparen} \setcounter{enumi}{4}
		\item $\lie{g}_2 = \lie{h}_2$.
	\end{enumparen}
\end{definition}

We shall show that, if $(\lie{g}, \lie{h}, \orbit{O})$ is a minimal NDD tuple with non-abelian $\lie{h}$, then $(\lie{g}, \lie{h})$ is of type $A_1$ or $BC'_1$.
We divide the proof into several steps.

Let $(\lie{g}, \lie{h}, \orbit{O})$ be a minimal NDD tuple with non-abelian $\lie{h}$.
By Proposition \ref{prop:BasicMinimal}, $\lie{h}$ is simple and has real rank one.

\begin{lemma} \label{lem:ReductionLemmaStep1}
	There exists an element $h \in \lie{h}^{-\theta}$ satisfying the following conditions.
	\begin{enumparen}
		\item $\lie{g} = \bigoplus_{i \in \ZZ} \lie{g}_i$.
		\item $\lie{h}_1 \neq 0$.
		\item For a positive integer $c$, one has $\lie{g}_c\cap \overline{\orbit{O}} \cap \lie{h}^\perp \neq \set{0}$ if and only if $c = 1$.
	\end{enumparen}
\end{lemma}

\begin{proof}
	Take a non-zero element $h \in \lie{a}_H$ such that $\lie{h} = \bigoplus_{-2\leq i \leq 2} \lie{h}_i$ and $\lie{h}_1 \neq 0$.
	By the representation theory of $\lieSL(2,\RR)$, every eigenvalue of $\ad_{\lie{g}}(h)$ is in $\frac{1}{2}\ZZ$,
	and hence we have $\lie{g} = \bigoplus_{i \in \ZZ} \lie{g}_{i/2}$.
	By Lemma \ref{lem:ExistenceWeightVector}, there exists a vector $X \in \overline{\orbit{O}} \cap \lie{h}^\perp \cap \lie{g}_c$ for some positive $c \in \frac{1}{2}\ZZ$.
	We take $X$ so that $c$ is maximal.

	If $\lie{h}_c = 0$, then $(\lie{g}, \lie{h}, \orbit{O})$ is not minimal by Corollary \ref{cor:SufficientNonRoot}.
	Hence we may assume that $\lie{h}_c \neq 0$, and then $c$ is either $1$ or $2$.
	By Corollary \ref{cor:MinimalGeneratedBy}, we have $\lie{g} = \bigoplus_{i \in \ZZ} \lie{g}_{ci}$.
	If $c = 2$, we have $\lie{g}_1 = 0$.
	This contradicts $\lie{h}_1 \neq 0$, and hence $c = 1$.
	Therefore we have shown the assertion.
\end{proof}

Fix an element $h \in \lie{h}^{-\theta}$ satisfying the conditions in Lemma \ref{lem:ReductionLemmaStep1}.

\begin{lemma} \label{lem:ReductionLemmaStep2}
	For any $d \geq 2$, if $\lie{h}_d = 0$, then one has $\lie{g}_d = 0$.
\end{lemma}

\begin{proof}
	Take a non-zero element $X$ in $\lie{g}_1\cap \overline{\orbit{O}} \cap \lie{h}^\perp$ by Lemma \ref{lem:ReductionLemmaStep1} (3).
	Let $d$ be the maximal positive integer such that $\lie{g}_d \neq 0$ and $p\colon \lie{g}\rightarrow \lie{g}_d$ be the projection with respect to the direct sum decomposition $\lie{g} = \bigoplus_{i \in \ZZ} \lie{g}_i$.
	Since $X$ does not belong to any proper ideal of $\lie{g}$ by Proposition \ref{prop:BasicMinimal} (6), $\Ad(G)X$ spans $\lie{g}$.
	Hence there exists $g \in G$ such that $p(\Ad(g)X) \neq 0$.
	Then we have
	\begin{align*}
		\overline{\orbit{O}} \ni \lim_{t \to \infty} e^{-dt} \Ad(\exp(th))\Ad(g)X = p(\Ad(g)X) \in \lie{g}_d.
	\end{align*}
	This shows that $\lie{g}_d \cap \overline{\orbit{O}} \neq \set{0}$.
	
	If $\lie{h}_d = 0$, then we have $\lie{g}_d \subset \lie{h}^\perp$, which contradicts the minimality of $(\lie{g}, \lie{h}, \orbit{O})$ by Corollary \ref{cor:SufficientNonRoot}.
	Since $\lie{h} = \lie{h}_{-2} \oplus \lie{h}_{-1} \oplus \lie{h}_0 \oplus \lie{h}_1 \oplus \lie{h}_2$, the assertion follows.
\end{proof}

\begin{lemma} \label{lem:ReductionLemmaStep3}
	Suppose that $\lie{g}_2 \neq 0$.
	Let $X$ be a non-zero element in $\overline{\orbit{O}} \cap \lie{g}_1$.
	Then one has
	\begin{align*}
		\lie{g}_2 = [X, \lie{g}_1] \oplus \lie{g}^{\theta(X)}_2 =  \lie{h}_2 \oplus \lie{g}^{\theta(X)}_2.
	\end{align*}
\end{lemma}

\begin{proof}
	Note that $[X, \lie{g}]^\perp = \lie{g}^X$.
	This implies that the orthogonal complement of $[X, \lie{g}_1]$ in $\lie{g}_2$ with respect to the inner product $-(\cdot, \theta(\cdot))$ is $\lie{g}_2^{\theta(X)}$, and we have $[X, \lie{g}_1] \oplus \lie{g}_2^{\theta(X)} = \lie{g}_2$.
	Hence it suffices to show that $\dim([X, \lie{g}_1]) \leq \dim(\lie{h}_2)$ and $\lie{h}_2 \cap \lie{g}_2^{\theta(X)} = 0$.

	For any $Y \in \lie{g}_1$, we have
	\begin{align*}
		\overline{\orbit{O}} &\ni \lim_{t \to \infty} e^{-2t} \Ad(\exp(th))\Ad(\exp(Y))X \\
		&= \lim_{t \to \infty} e^{-2t} \Ad(\exp(th))(X + [Y, X]) = [Y, X].
	\end{align*}
	This implies that $[X, \lie{g}_1] \subset \overline{\orbit{O}}$.
	If $[X, \lie{g}_1] \cap \lie{h}^\perp \neq 0$, then we have $\overline{\orbit{O}}\cap \lie{g}_2 \cap \lie{h}^\perp \neq \set{0}$, which contradicts Lemma \ref{lem:ReductionLemmaStep1} (3).
	Hence we have $[X, \lie{g}_1] \cap \lie{h}^\perp = 0$ and $\dim([X, \lie{g}_1]) \leq \dim(\lie{h}_2)$.

	Since $\lie{h}$ is a simple Lie algebra of real rank one, for any non-zero $x \in \lie{h}_2$, the subalgebra generated by $x$ and $\theta(x)$ is isomorphic to $\lieSL(2,\RR)$.
	This implies that $[\theta(X), x] \neq 0$ and hence $\lie{h}_2 \cap \lie{g}_2^{\theta(X)} = 0$.
	We have shown the lemma.
\end{proof}

\begin{corollary} \label{cor:ReductionLemmaStep3}
	One has $\overline{\orbit{O}} \cap \lie{g}_1 \cap \lie{h}^\perp\subset \lie{g}^{\lie{h}_1}$.
\end{corollary}

\begin{proof}
	If $\lie{g}_2 = 0$, then the assertion is trivial.
	Assume that $\lie{g}_2 \neq 0$.
	In the proof of Lemma \ref{lem:ReductionLemmaStep3}, we have shown that $[X, \lie{g}_1] \cap \lie{h}^\perp = 0$ for any $X \in \overline{\orbit{O}} \cap \lie{g}_1$.
	This implies that $[X, \lie{h}_1] = 0$ for any $X \in \overline{\orbit{O}} \cap \lie{g}_1 \cap \lie{h}^\perp$.
\end{proof}

\begin{lemma} \label{lem:ReductionLemmaStep4}
	If $\lie{g}_2 \neq 0$, then $\lie{g}$ is simple.
\end{lemma}

\begin{proof}
	Assume that $\lie{g}$ is not simple and $\lie{g} = \lie{g}' \oplus \lie{g}''$ is a direct sum of two proper non-zero ideals.
	Note that $\overline{\orbit{O}}$ is written as $\overline{\orbit{O}} = \overline{\orbit{O}'} + \overline{\orbit{O}''}$ for some nilpotent orbits $\orbit{O}'$ and $\orbit{O}''$ in $\lie{g}'$ and $\lie{g}''$, respectively.
	Take a non-zero element $X \in \overline{\orbit{O}} \cap \lie{g}_1 \cap \lie{h}^\perp$ and write $X = X' + X''$ with $X' \in \lie{g}'$ and $X'' \in \lie{g}''$.
	Then we have $X', X'' \in \overline{\orbit{O}}$.

	By Proposition \ref{prop:BasicMinimal} (6), $X'$ and $X''$ are non-zero.
	By Lemma \ref{lem:ReductionLemmaStep3}, we have
	\begin{align*}
		\dim([X, \lie{g}_1]) = \dim([X', \lie{g}'_1]) = \dim([X'', \lie{g}''_1]) = \dim(\lie{h}_2).
	\end{align*}
	On the other hand, we have $[X, \lie{g}_1] = [X', \lie{g}'_1] + [X'', \lie{g}''_1]$.
	This contradicts the above equality, and we have shown the lemma.
\end{proof}

Summarizing Propositions \ref{prop:BasicMinimal} and \ref{prop:BasicMinimalH}, and Lemmas \ref{lem:ReductionLemmaStep1}, \ref{lem:ReductionLemmaStep2}, \ref{lem:ReductionLemmaStep3} and \ref{lem:ReductionLemmaStep4}, we have the following theorem.

\begin{theorem} \label{thm:ReductionLemma}
	If $(\lie{g}, \lie{h}, \orbit{O})$ is a minimal NDD tuple with non-abelian $\lie{h}$, then $(\lie{g}, \lie{h})$ is of type $A_1$ or $BC'_1$.
\end{theorem}

\subsection{Reduction by a \texorpdfstring{$\lie{q}$}{q}-orthogonal system} \label{subsect:ReductionStronglyOrthogonal}

The proof of Theorem \ref{thm:MainTargetMinimal} relies on classification results and certain explicit computations.
In this subsection, we prove a proposition that helps reduce the amount of computation.

Let $(\lie{g}, \lie{h}, \orbit{O})$ be an NDD tuple of type $A_1$ or $BC'_1$.
Fix a non-zero element $h \in \lie{h}^{-\theta}$ as in Definitions \ref{def:typeA} and \ref{def:typeBC}.
Set $\lie{a}_H \coloneq \RR h$.
Then we have a grading $\lie{g} = \bigoplus_{i=-2}^2 \lie{g}_i$ associated with $h$.
Note that $\lie{g}_2$ is zero in the case of type $A_1$.

Let $\sigma$ be the involution on $\lie{g}$ such that $\lie{g}^\sigma = \lie{g}_{-2} \oplus \lie{g}_0 \oplus \lie{g}_2$.
Fix a maximal abelian subspace $\lie{a}$ of $\lie{g}^{-\theta}$ containing $\lie{a}_H$.
Then $\lie{a} \subset \lie{g}^{-\theta}_0$ and $\lie{g}_1$ is $\lie{a}$-stable.
Take roots $\set{\lambda_1, \ldots, \lambda_r}$ in $\Delta(\lie{g}_1, \lie{a})$
and root vectors $x_i \in \lie{g}_{\lambda_i}$ such that
\begin{align}
	[x_i, x_j] = 0, \quad [x_i, \theta(x_j)] = 0 \quad (\forall i \neq j), \label{eqn:QOrthogonal}
\end{align}
and $r$ is as large as possible.
Note that $\set{\lambda_1, \ldots, \lambda_r}$ may not be a set of strongly orthogonal roots in general.
The set $\set{x_1, \ldots, x_r}$ is called a $\lie{q}$-orthogonal system in \cite{Ma79_orbits}.

\begin{fact}[T.\ Matsuki {\cite[Theorem 2]{Ma79_orbits}}] \label{fact:Matsuki}
	$\bigoplus_{i=1}^r \RR (x_i - \theta(x_i))$ is a maximal abelian subspace in $\lie{g}^{-\sigma, -\theta} = (\lie{g}_{-1}\oplus \lie{g}_1) \cap \lie{g}^{-\theta}$.
\end{fact}

We shall construct a suitable candidate for the subalgebra $\lie{g}'$ of $\lie{g}$ in Proposition \ref{prop:NonMinimalSubalgebra} using the $\lie{q}$-orthogonal system $\set{x_1, \ldots, x_r}$.
Set $\lie{c} \coloneq \bigoplus_{i=1}^r \RR x_i$.
Let $\lie{g'}$ be the subalgebra of $\lie{g}$ generated by $\lie{c}$, $\theta(\lie{c})$ and $\lie{a}$.

\begin{lemma} \label{lem:ReductionStronglyOrthogonalBasic}
	The subalgebra $\lie{g'}$ satisfies the following conditions.
	\begin{enumparen}
		\item $\lie{g'}$ is $\theta$-stable.
		\item $\lie{g'} = \lie{a} \oplus \lie{c} \oplus \theta(\lie{c})$.
		\item $\lie{g'}$ is isomorphic to a direct sum of some copies of $\lieSL(2,\RR)$ and an abelian ideal.
	\end{enumparen}
\end{lemma}

\begin{proof}
	The assertions are clear from the definition of $\lie{g}'$ and the relations \eqref{eqn:QOrthogonal}.
\end{proof}

\begin{lemma} \label{lem:StronglyOrthogonalHCompatible}
	Assume that $\lie{c} = (\lie{c} \cap \lie{h}) \oplus (\lie{c} \cap \lie{h}^\perp)$.
	Then one has $\lie{g'} = (\lie{g'} \cap \lie{h}) \oplus (\lie{g'} \cap \lie{h}^\perp)$.
\end{lemma}

\begin{proof}
	By Lemma \ref{lem:ReductionStronglyOrthogonalBasic} (2), we have
	\begin{align*}
		\lie{g'} &= \lie{a} \oplus \lie{c} \oplus \theta(\lie{c}), \\
		\lie{g'} \cap \lie{h} &\supset \lie{a}_H \oplus (\lie{c} \cap \lie{h}) \oplus (\theta(\lie{c}) \cap \lie{h}), \\
		\lie{g'} \cap \lie{h}^\perp &\supset (\lie{a}_H^\perp \cap \lie{a}) \oplus(\lie{c} \cap \lie{h}^\perp) \oplus (\theta(\lie{c}) \cap \lie{h}^\perp).
	\end{align*}
	From this and the assumption $\lie{c} = (\lie{c} \cap \lie{h}) \oplus (\lie{c} \cap \lie{h}^\perp)$, we obtain $\lie{g'} = (\lie{g'} \cap \lie{h}) \oplus (\lie{g'} \cap \lie{h}^\perp)$.
\end{proof}

Hereafter, assume that $\lie{g}$ has a simple ideal not isomorphic to $\lieSL(2,\RR)$.
This assumption is satisfied if $\lie{h}$ is not isomorphic to $\lieSL(2,\RR)$ since $\lie{h}$ is not contained in any proper ideal of $\lie{g}$.
See Definition \ref{def:typeA} (4).

\begin{lemma} \label{lem:ReductionRealStronglyOrthogonal}
	Assume that $\lie{c} \cap \lie{h}^\perp \cap \overline{\orbit{O}} \neq \set{0}$ and $\lie{c} = (\lie{c} \cap \lie{h}) \oplus (\lie{c} \cap \lie{h}^\perp)$.
	Then $(\lie{g}, \lie{h}, \orbit{O})$ is not minimal.
\end{lemma}

\begin{proof}
	By the assumption that $\lie{g}$ has a simple factor not isomorphic to $\lieSL(2,\RR)$, we have $\lie{g} \neq \lie{g'}$.
	By Lemma \ref{lem:StronglyOrthogonalHCompatible}, we have $\lie{g'} = (\lie{g'} \cap \lie{h}) \oplus (\lie{g'} \cap \lie{h}^\perp)$.
	By definition, any non-zero vector in $\lie{c} \subset \lie{g}_1$ is an $\lie{a}_H$-weight vector.
	By Proposition \ref{prop:NonMinimalSubalgebra}, $(\lie{g}, \lie{h}, \orbit{O})$ is not minimal.
\end{proof}

A direct verification of the conditions of Lemma \ref{lem:ReductionRealStronglyOrthogonal} still requires explicit computations.
The following proposition has a narrower range of applicability, but it allows us to check the conditions of Lemma \ref{lem:ReductionRealStronglyOrthogonal} using only routine computations.

\begin{proposition} \label{prop:ReductionRealStronglyOrthogonal}
	Assume the following conditions.
	\begin{enumparen}
		\item There exist two $\theta$-stable ideals $\lie{i}$ and $\lie{j}$ of $\lie{g}^{\theta\sigma}$ such that
		\begin{align*}
			&\lie{g}^{\theta\sigma} = \lie{i} \oplus \lie{j}, \\
			&\lie{g}^{-\sigma, -\theta} \cap \lie{h} = \lie{i}^{-\theta}, \\
			&\lie{g}^{-\sigma, -\theta} \cap \lie{h}^\perp = \lie{j}^{-\theta}.
		\end{align*}
		\item The restriction to $\lie{g}_0\cap \lie{k}$ of the orthogonal projection from $\lie{g}$ to $\lie{j}\cap \lie{k}$ is surjective.
	\end{enumparen}
	Then $(\lie{g}, \lie{h}, \Ad(G)X)$ is not minimal for any $0\neq X \in \lie{g}_1 \cap \lie{h}^\perp$.
\end{proposition}

\begin{proof}
	We shall verify the two assumptions of Lemma \ref{lem:ReductionRealStronglyOrthogonal}.
	Define a map $\varphi \colon \lie{g}_1 \rightarrow \lie{g}^{-\sigma, -\theta} = (\lie{g}_{-1} \oplus \lie{g}_{1}) \cap \lie{g}^{-\theta}$ by $\varphi(X) = X - \theta(X)$.
	Then $\varphi$ is bijective and $(\lie{g}_0\cap \lie{k})$-equivariant, and we have
	\begin{align*}
		\varphi(\lie{g}_1 \cap \lie{h}) &= \lie{g}^{-\sigma, -\theta} \cap \lie{h} = \lie{i}^{-\theta}, \\
		\varphi(\lie{g}_1 \cap \lie{h}^\perp) &= \lie{g}^{-\sigma, -\theta} \cap \lie{h}^\perp = \lie{j}^{-\theta}.
	\end{align*}
	Since $\lie{g}^{\theta\sigma} = \lie{i} \oplus \lie{j}$, any maximal abelian subspace in $\lie{g}^{-\theta, -\sigma}$ is of the form $\lie{c}' \oplus \lie{c}''$ with $\lie{c}' \subset \lie{i}$ and $\lie{c}'' \subset \lie{j}$.
	Since $\varphi(\lie{c})$ is a maximal abelian subspace in $\lie{g}^{-\sigma, -\theta}$ by Fact \ref{fact:Matsuki}, we obtain 
	\begin{align*}
		\lie{c} = (\lie{c} \cap \lie{h}) \oplus (\lie{c} \cap \lie{h}^\perp).
	\end{align*}
	This is the second assumption in Lemma \ref{lem:ReductionRealStronglyOrthogonal}.

	Let $0\neq X \in \lie{g}_1 \cap \lie{h}^\perp$.
	Then we have $\varphi(X) \in \lie{j}^{-\theta}$.
	Let $K_0$ be the analytic subgroup of $G$ with the Lie algebra $\lie{g}_0\cap \lie{k}$.
	Since $\varphi(\lie{c}) \cap \lie{h}^\perp$ is a maximal abelian subspace in $\lie{j}^{-\theta}$, assumption (2) implies that there exists $k \in K_0$ such that $\Ad(k)\varphi(X) \in \varphi(\lie{c}) \cap \lie{h}^\perp$.
	Hence we obtain
	\begin{align*}
		0\neq \Ad(k)X \in \lie{c} \cap \lie{h}^\perp \cap \overline{\Ad(G)X}.
	\end{align*}
	This shows the first assumption in Lemma \ref{lem:ReductionRealStronglyOrthogonal}, that is, $\lie{c} \cap \lie{h}^\perp \cap \overline{\Ad(G)X} \neq \set{0}$.
	We have shown the proposition.
\end{proof}
